\documentclass{article}
\usepackage[UTF8]{ctex}
\usepackage{amsmath,amssymb,amsfonts}
\usepackage{enumitem}
\usepackage{geometry}
\geometry{a4paper, margin=1in}
% === 错题本排版优化（不参与题目标识） ===
\setlength{\parindent}{0pt}
\setlength{\parskip}{0.35em plus 0.10em minus 0.08em}
\setlength{\emergencystretch}{3em}
\setlist{itemsep=0.25em, topsep=0.40em, parsep=0pt}
\allowdisplaybreaks[2]
% === 错题本排版优化结束 ===
\begin{document}
\section*{西南大学2025年数学分析真题}
\section*{65．西南大学2025年数学分析试卷}

\subsection{题目1}
1． \(\lim _{n \rightarrow \infty}\left(\frac{1}{\ln (n+1)-\ln n}-n\right)\) ．

\subsection{题目2}
2． \(\lim _{x \rightarrow 0}\left(\frac{\cos x}{\cos 2 x}\right)^{\frac{1}{\ln \left(1+x^{2}\right)}}\) ．

\subsection{题目3}
3． \(\int_{-\pi}^{\pi}(\sin 2 x)\left(\cos \left(x^{2}\right)+e^{x}\right) \mathrm{d} x\) ．

\subsection{题目4}
1．设 \(F\) 有二阶连续偏导数，且 \(z=f(x, y)\) 是由 \(F(x, x+y, x+y+z)=0\) 所确定的隐函数，求 \(\frac{\partial^{2} z}{\partial x^{2}}\) ．

\subsection{题目5}
2．计算

\[
\iiint_{V}\left(x^{2}+y^{2}+z^{\frac{1}{2}}\right) \mathrm{d} x \mathrm{~d} y \mathrm{~d} z
\]

其中 \(V\) 是由 \(\left\{\begin{array}{l}z=y \\ z=0\end{array}\right.\) 绕 \(z\) 轴旋转一周的曲面与 \(z=1\) 所围成的区域．

\subsection{题目6}
3．求 \(f(x, y)=x^{2} y(4-x-y)\) 在由 \(x\) 轴，\(y\) 轴与 \(x+y=6\) 所围区域 \(D\) 中的最大值与最小值．

三．（15 分）设 \(f(x)\) 在 \((-\infty, a)\) 上连续，且 \(\lim _{x \rightarrow-\infty} f(x)\) 与 \(\lim _{x \rightarrow a^{-}} f(x)\) 存在．证明：\(f(x)\) 在 \((-\infty, a)\) 上一致连续。

四．（15 分）设 \(f(x)\) 在 \([a, b]\) 上连续，在 \((a, b)\) 上可导，且存在 \(\xi \in(a, b)\) ，满足 \(f^{\prime}(\xi)=0\) ．证明：存在\(\eta \in(a, b)\) ，使得 \(f^{\prime}(\eta)-f(\eta)+f(a)=0\) 。

五．（15 分）设 \(f(x)\) 在 \([a, b]\) 上连续，且对任意的 \(x \in(a, b)\) ，有 \(f(x)>f(a)=f(b)\) ．证明：存在\(\xi \in(a, b)\) ，使得

\[
\frac{\int_{a}^{\xi} f(x) \mathrm{d} x}{\xi-a}=\frac{\int_{\xi}^{b} f(x) \mathrm{d} x}{b-\xi}
\]

六．（15 分）设 \(f(x)=\sum_{n=2}^{\infty} \frac{1}{n^{x}}, x \in(1,+\infty)\) ，常数 \(a>1\) ．判断 \(\int_{a}^{+\infty} f(x) \mathrm{d} x\) 的敛散性并证明．

七．（15 分）设 \(f(x)\) 在 \([a, b]\) 上有连续的导函数，记 \(J=\frac{1}{b-a} \int_{a}^{b} f(x) \mathrm{d} x\) ．证明：

\[
\frac{1}{b-a} \int_{a}^{b}|f(x)-J| \mathrm{d} x \leq \int_{a}^{b}\left|f^{\prime}(x)\right| \mathrm{d} x .
\]

八．（15 分）证明：含参量积分 \(\int_{0}^{+\infty} \frac{y \cos x y}{1+y^{2}} \mathrm{~d} y\) 在 \((0,+\infty)\) 上不一致收敛．
\end{document}
